\section{Les interprétations idéalistes et compatibles}
\subsection{Interprétation statistique (sans modification des lois quantiques)}
Le module du coefficient en facteur d’un vecteur d’état, au sein d’une combinaison linéaire
d’états, peut être décrit exclusivement comme la racine carrée de la probabilité normalisée
[= nombre de cas / nombre total de mesures] de mesurer la valeur propre associée (d’une
observable donnée), soit, une fois multipliée par le nombre de mesures effectuées dans un
laboratoire (sur un type de système préparé dans des conditions identiques), le nombre de
fois où cette valeur doit être obtenue comme résultat. Ainsi, le sens donné à une
superposition quantique peut être de nature purement statistique et caractériser un ensemble
de systèmes étudiés, au lieu d’un système unique. Le paradoxe de la superposition d’états
disparaît donc de par son application conceptuelle à un ensemble de systèmes. A.Einstein,
qui était un partisan de cette vision « ensembliste », écrit dans une lettre à E.Schrödinger : «
la fonction $\Psi$ ne décrit pas l’état d’un système unique mais (statistiquement) un ensemble de
systèmes ». Vision qui ne spécifie aucunement la nature de la réalité d’un système
antérieurement à la mesure (empirique) et peut donc tout aussi bien se classer dans le
réalisme que dans les diverses variantes de l’idéalisme.
\subsection{Point de vue informationnel pur (sans modification des lois
 quantiques)}
 Le vecteur d’état peut être vu comme contenant, donnant une information accessible sur le
système et non un élément de réalité caractérisant ses propriétés [impulsion, spin...]
(comme le ‘Quantum Bayesianism (QBism)’ de Fuchs, Schack, Mermin). La superposition
quantique est bien possible mais n’induit alors plus de conflit logique. L’aspect probabiliste
de la théorie ne fait plus que quantifier et paramétrer une ignorance, sans révéler un
comportement aléatoire (ou autre) intrinsèque au système. La réduction du vecteur d’état
n’est plus qu’un affûtage de la connaissance. Néanmoins, le vecteur d’état, quoique
‘mental’, continue à concerner un système physique unique. Ce point de vue laisse ouverte
la question ontologique sur la réalité de l’objet, sans trancher entre le réalisme et
l’idéalisme. Faut-il abandonner tout espoir que les sciences physiques nous apprennent
quelque chose sur l’essence du monde réel ? Dans ces écrits des premières heures,
W.Heisenberg adoptait une approche assez proche du point de vue purement informationnel,
avant, plus tard, de nuancer ses propos.
\subsection{Interprétation des corrélations (sans modification des lois
 quantiques)}
 La Mécanique quantique est considérée ici comme une description théorique de nature
stochastique (caractère aléatoire) fournissant en effet des probabilités associées à des
résultats de mesure, mais il n’y a pas de réalité (existence) du système avant sa mesure, y
compris lors d’une série de mesures intermédiaires – aux résultats corrélés formellement par
les états successivement projetés du système étudié [état final d’une mesure = état initial de
la suivante]. Par exemple, si un photon est capté après son passage dans une fibre optique,
son parcours ne recouvre pas de réalité physique (ou bien représente une question superflue
dans une version moins drastique de cette interprétation ressemblant alors davantage à un
modèle effectif). Il n’est alors plus utile de se questionner physiquement sur le concept de
superposition d’états ni donc sur le postulat de réduction de fonction d’onde (‘saut
discontinu’), notions qui deviennent de simples artefacts du formalisme. Dans le sens où
n’existe que ce qui est observé dans le monde empirique, cette interprétation, dite des
corrélations, correspond à la vision hautement idéaliste de J.G.Fichte se passant tout
simplement du « noumène ».
